%% Appendix: Complementary Quantitative Verification — Form-Only Respondents.
%% Aggregate analysis of seven respondents who returned the structured questionnaire
%% after reading the pre-read paper independently, without participating in a
%% semi-structured interview. Distinct instrument from the interview series.

This appendix reports the aggregate quantitative findings of a complementary verification instrument: seven respondents who returned the structured questionnaire described in Appendix~\ref{ape:methodology} after reading the pre-read paper independently, without participating in the semi-structured interview. The instrument is distinct from the fifteen-session interview series reported in Appendix~\ref{ape:interviews}: there is no qualitative transcript to triangulate against, no per-respondent analytical coding, and no participation in the saturation analysis. The role of this complementary instrument is to provide an independent quantitative check on the framework's reception, gathered from practitioners who engaged with the paper without the conversational anchoring of an interview.

\section*{Instrument and respondent attribution}

The seven respondents returned the structured questionnaire between 8 June 2026 and 10 June 2026. Three of the seven respondents share the email domain of the public-sector employer that supplied three of the interview series participants (Interviews 13, 14, and 15), suggesting they belong to the same organizational context; the remaining four respondents declared no traceable affiliation through the email surface. All seven are treated as anonymous quantitative inputs in this appendix; no per-respondent profile is reported because none was collected during the questionnaire.

The complementary instrument is reported separately rather than aggregated with the thirteen interview-and-questionnaire responses to preserve the methodological distinction between the two verification surfaces: the interview series provides qualitative depth with per-session triangulation, while the form-only instrument provides quantitative breadth with no qualitative anchor. The aggregate counts that follow report N equal to seven for the form-only group and N equal to thirteen for the interview-and-questionnaire group (the two interview participants whose questionnaire is still pending are not included in the aggregate analysis).

\section*{Aggregate findings}

The headline of the comparison is convergence on the framework's structure with three substantive divergences that independently reinforce qualitative findings from the interview series. The convergent items are reported first; the divergent items are reported as findings in their own right in the subsection that follows.

\subsection*{Section~1 general framing}

\begin{longtable}{p{6cm} c c c}
\toprule
\textbf{Item} & \textbf{Interview+form} & \textbf{Form-only} & \textbf{Delta} \\
& mean (sd) & mean (sd) & \\
\midrule
\endhead
Four dimensions reflect practice & 4.38 (0.65) & 4.14 (0.38) & $-0.24$ \\
Modular monolith as deliberate destination & 4.62 (0.65) & 4.14 (1.07) & $-0.48$ \\
G1--G6 ordering matches team approach & 4.08 (0.95) & 4.00 (0.82) & $-0.08$ \\
Set complete (no missing guideline) & 4.54 (0.52) & 4.29 (1.11) & $-0.25$ \\
\bottomrule
\end{longtable}

\noindent
The four general-framing items return aggregate means within plus or minus zero point five between the two groups, indicating strong convergence on the dissertation's framework structure. The form-only respondents endorse the four-dimension framing, the modular monolith as a deliberate destination, the G1--G6 ordering, and the completeness of the six-guideline set at substantively equivalent levels to the interview-and-questionnaire respondents.

\subsection*{Per-guideline ratings}

\begin{longtable}{p{2cm} c c c c c}
\toprule
\textbf{Guideline} & \textbf{Applic.} & \textbf{Clarity} & \textbf{Importance} & \textbf{Adoption} & \textbf{Completeness} \\
\midrule
\endhead
G1 IF mean & 4.23 & 4.31 & 4.77 & 4.00 & 4.31 \\
G1 FO mean & 4.43 & 3.86 & 4.71 & 3.86 & 4.14 \\
G1 delta & $+0.20$ & $-0.45$ & $-0.06$ & $-0.14$ & $-0.17$ \\
\midrule
G2 IF mean & 4.31 & 3.92 & 4.54 & 4.23 & 4.15 \\
G2 FO mean & 4.29 & 4.14 & 4.57 & 4.00 & 3.86 \\
G2 delta & $-0.02$ & $+0.22$ & $+0.03$ & $-0.23$ & $-0.29$ \\
\midrule
G3 IF mean & 4.23 & 4.31 & 4.31 & 4.00 & 4.23 \\
G3 FO mean & 4.14 & 4.00 & 4.43 & 3.86 & 3.86 \\
G3 delta & $-0.09$ & $-0.31$ & $+0.12$ & $-0.14$ & $-0.37$ \\
\midrule
G4 IF mean & 4.31 & 4.15 & 4.38 & 3.85 & 4.15 \\
G4 FO mean & 3.86 & 3.86 & 4.14 & 3.43 & 3.86 \\
G4 delta & $-0.45$ & $-0.29$ & $-0.24$ & $-0.42$ & $-0.29$ \\
\midrule
G5 IF mean & 4.31 & 4.46 & 4.46 & 4.38 & 4.08 \\
G5 FO mean & 4.29 & 4.14 & 4.29 & 4.00 & 4.14 \\
G5 delta & $-0.02$ & $-0.32$ & $-0.17$ & $-0.38$ & $+0.06$ \\
\midrule
G6 IF mean & 4.69 & 4.62 & 4.77 & 4.54 & 4.54 \\
G6 FO mean & 4.57 & 4.43 & 4.71 & 4.00 & 4.57 \\
G6 delta & $-0.12$ & $-0.19$ & $-0.06$ & $-0.54$ & $+0.03$ \\
\bottomrule
\end{longtable}

\noindent
Twenty-six of the thirty per-guideline aggregate means are within plus or minus zero point four between the two groups, indicating convergence on most of the per-guideline judgment surface. The four deltas that exceed plus or minus zero point four are concentrated on the adoption dimension and on G4: the G6 adoption likelihood drops by zero point five four between the interview-and-questionnaire group and the form-only group, the G4 applicability and adoption drop by zero point four five and zero point four two respectively, and the G1 clarity drops by zero point four five.

\subsection*{Named gap recognition}

\begin{longtable}{p{5cm} c c}
\toprule
\textbf{Named gap} & \textbf{IF count / N} & \textbf{FO count / N} \\
\midrule
\endhead
The Enforcement Gap (G1) & 13 / 13 & 5 / 7 \\
The Incremental Maintainability Degradation (G2) & 12 / 13 & 5 / 7 \\
The Progressive Scalability Gap (G3) & 10 / 13 & 6 / 7 \\
The Migration Readiness Gap (G4) & 11 / 13 & 5 / 7 \\
The Deployment-Architecture Mismatch (G5) & 7 / 13 & 3 / 7 \\
The Observability Deficit (G6) & 12 / 13 & 4 / 7 \\
\bottomrule
\end{longtable}

\noindent
The two groups recognize the named gaps at substantively comparable proportional rates. The Enforcement Gap, the Incremental Maintainability Degradation, and the Progressive Scalability Gap are the most consistently recognized in both groups. The Deployment-Architecture Mismatch is the least consistently recognized in both groups, converging with the dissertation's framing of G5 as situationally activated.

\subsection*{$\varepsilon_m$ calibration distribution}

\begin{longtable}{p{6cm} c c}
\toprule
\textbf{Response} & \textbf{IF count / N} & \textbf{FO count / N} \\
\midrule
\endhead
About right & 8 / 13 & 2 / 7 \\
Too strict to be practically achievable & 2 / 13 & 3 / 7 \\
Not confident enough to answer & 3 / 13 & 2 / 7 \\
\bottomrule
\end{longtable}

\noindent
This is the strongest distinct pattern between the two groups. The interview-and-questionnaire respondents endorse the $\varepsilon_m$ calibration as \emph{About right} in eight of the thirteen responses; the form-only respondents shift the distribution toward \emph{Too strict to be practically achievable}. The shift is the quantitative complement of the soft-gates-to-hard-gates reformulation proposed in Interview 12 (Renan) and is reported as a finding in the divergence subsection below.

\subsection*{Spectrum and gate evaluation}

\begin{longtable}{p{6cm} c c c}
\toprule
\textbf{Item} & \textbf{Interview+form} & \textbf{Form-only} & \textbf{Delta} \\
& mean (sd) & mean (sd) & \\
\midrule
\endhead
G3 progressive scalability spectrum (D0--D2) & 4.69 (0.63) & 4.71 (0.49) & $+0.02$ \\
G5 deployment spectrum (D0--D2) & 4.15 (0.80) & 4.29 (0.76) & $+0.13$ \\
Composite binary gates as enforcement mechanism & 3.77 (0.93) & 3.57 (0.53) & $-0.20$ \\
\bottomrule
\end{longtable}

\noindent
The three structural-mechanism Likerts return essentially equivalent aggregate means between the two groups, indicating convergent endorsement of the dissertation's spectrum framings and of the composite binary gates as a useful mechanism (with the qualification on the $\varepsilon_m$ calibration noted above).

\subsection*{Forced prioritization within G1--G6}

\begin{longtable}{p{3cm} c c c}
\toprule
\textbf{Guideline} & \textbf{IF mean rank} & \textbf{FO mean rank} & \textbf{Delta} \\
\midrule
\endhead
G1 Enforce Modular Boundaries & 2.62 & 2.71 & $+0.09$ \\
G2 Embed Maintainability & 3.23 & 3.14 & $-0.09$ \\
G3 Design for Progressive Scalability & 3.38 & 3.57 & $+0.19$ \\
G4 Promote Migration Readiness & 4.23 & 4.29 & $+0.06$ \\
G5 Streamline Deployment Strategy & 3.62 & 4.43 & $+0.81$ \\
G6 Introduce Observability Patterns & 3.15 & 4.00 & $+0.85$ \\
\bottomrule
\end{longtable}

\noindent
The G1, G2, G3, and G4 mean ranks return essentially identical between the two groups (within plus or minus zero point two), indicating convergent prioritization on four of the six guidelines. The two divergences are concentrated on G5 and G6: the form-only respondents place both at lower priority positions than the interview-and-questionnaire respondents, with the G6 delta of plus zero point eight five as the largest in the table.

\subsection*{Forced prioritization within G7--G12}

\begin{longtable}{p{3cm} c c c}
\toprule
\textbf{Guideline} & \textbf{IF mean rank} & \textbf{FO mean rank} & \textbf{Delta} \\
\midrule
\endhead
G7 Team-Architecture Balance & 3.31 & 2.71 & $-0.60$ \\
G8 SRE and DevOps Maturity & 3.69 & 3.43 & $-0.26$ \\
G9 Frictionless Onboarding & 3.77 & 5.14 & $+1.37$ \\
G10 Actionable Design Patterns & 2.92 & 2.71 & $-0.21$ \\
G11 Contextual Adaptation & 3.62 & 3.86 & $+0.24$ \\
G12 Trade-offs over Dogma & 3.46 & 5.00 & $+1.54$ \\
\bottomrule
\end{longtable}

\noindent
G10 Actionable Design Patterns occupies the top position in both groups, converging on the practitioner expectation that the deferred dimensions deliver value through concrete patterns. The form-only group elevates G7 Team-Architecture Balance (delta of minus zero point six) and demotes G9 Frictionless Onboarding (delta of plus one point three seven) and G12 Trade-offs over Dogma (delta of plus one point five four) relative to the interview-and-questionnaire group. The pattern is the quantitative complement of the Organizational Alignment cross-axis reformulation that surfaced in the interview series.

\section*{Three divergent findings}

The convergence reported above strengthens the case for the framework's structural choices. The three divergences are reported below as substantive findings in their own right, each of which converges with a qualitative theme from the interview series.

\paragraph*{D1. G6 adoption likelihood drops in the form-only group.} The G6 adoption rating drops by zero point five four (from 4.54 to 4.00) between the interview-and-questionnaire and the form-only groups, and the G6 mean rank drops by zero point eight five between the two groups. The pattern converges with the timing-based and cost-conditioned critiques of G6 surfaced in Interviews 12, 13, and 14: G6 is endorsed as important in principle (importance rating remains at 4.71 in the form-only group) but its adoption is treated as conditional on team maturity, criticality, and budget. The dissertation records this convergence as additional empirical support for the cost-conditioned reading of G6 that the interview series produced.

\paragraph*{D2. $\varepsilon_m$ calibration is perceived as too strict by the form-only group.} The form-only group shifts the $\varepsilon_m$ calibration distribution toward \emph{Too strict to be practically achievable} (three of seven respondents, compared with two of thirteen in the interview-and-questionnaire group). The pattern converges directly with the soft-gates-to-hard-gates reformulation proposed in Interview 12 (Renan), which the dissertation records as a post-defense iteration candidate. The form-only data strengthens the case that the soft-gates framing is not idiosyncratic to a single voice but is shared by a substantive minority of independent readers.

\paragraph*{D3. G7 Team-Architecture Balance rises in the form-only prioritization.} The G7 mean rank in the deferred-dimensions ordering drops by zero point six between the two groups (from 3.31 to 2.71), making G7 the second-highest priority in the form-only group after G10 Actionable Design Patterns. The pattern converges with the Organizational Alignment cross-axis reformulation that surfaced in Interviews 9, 10, and 13, in which participants proposed that organizational concerns traverse the framework rather than being deferred to a parallel research direction. The form-only data adds independent quantitative weight to this reformulation as a research direction.

\section*{Open reflection preserved verbatim}

Six of the seven form-only respondents returned both optional reflection fields as no-content entries. The seventh returned the reflection on points not discussed (the questionnaire field corresponding to interview supplementation) with the following text, preserved verbatim because it expresses a distinct architectural framing not surfaced in prior sessions:

\begin{quote}
\emph{``No meu entendimento monolitos e microservi\c{c}os s\~ao dois extremos quando falamos de arquitetura. Nas experi\^encias que j\'a tivemos consumindo microservi\c{c}os notamos um excesso de requisi\c{c}\~oes e tratativas para recuperar informa\c{c}\~oes simples, mas com um m\'inimo de relacionamentos. Dessa forma optamos por servi\c{c}os um pouco mais complexos que respodessem diretamente \`a necessidade do solicitante. Seria um meio termo em microservi\c{c}os e um backend puro.''} (Form-only respondent, 10 June 2026)

\emph{[English: ``In my understanding monoliths and microservices are two extremes when we speak of architecture. In the experiences we have had consuming microservices we noticed excessive requests and processing to retrieve simple information with minimal relationships. We therefore opted for slightly more complex services that respond directly to the requester's need. This would be a middle ground between microservices and a pure backend.'']}
\end{quote}

\noindent
The reflection articulates the same target-architecture positioning the dissertation defends through the modular monolith framing: a deliberate middle ground between two extremes, chosen for the operational consequences each extreme would impose. The convergence is recorded here because it expresses the positioning in a vocabulary independent of the dissertation's own (the respondent describes the middle ground in terms of request granularity rather than in terms of module boundaries, which is itself a finding worth preserving).

\section*{Limitations of the complementary instrument}

\paragraph*{L1. No qualitative triangulation surface.} The form-only instrument provides numerical responses but no qualitative transcript against which to verify, code, or interpret the ratings. The convergence and divergence reported above is therefore aggregate quantitative signal only, with the depth of the per-session analysis available only for the interview series.

\paragraph*{L2. Sample size of seven.} The aggregate means and the divergence deltas are computed over a sample of seven respondents. The means are reported with sample standard deviations, but no inferential statistics are claimed; the role of the instrument is convergence checking, not hypothesis testing.

\paragraph*{L3. Self-selection effect distinct from the interview series.} The form-only respondents engaged with the paper without the conversational anchoring of an interview. This is a feature rather than a defect for convergence-checking purposes, since it removes the interview-induced anchoring as a potential confound, but it also means the form-only group may be skewed toward respondents who read the paper without engaging with the researcher.

\paragraph*{L4. Ranking grids returned with ties in some responses.} Some form-only respondents used only some of the available rank positions, returning ties rather than strict ordinal rankings. The mean ranks reported above include all returned positions; the threats-to-validity discussion of Chapter~5 documents this measurement-level limitation for both groups.

\paragraph*{L5. Open reflections almost entirely absent.} Six of the seven form-only respondents returned both optional reflection fields blank. The single verbatim reflection preserved above is the entirety of the qualitative material the complementary instrument produced.

\section*{Contributions to the conclusions chapter}

\begin{enumerate}[itemsep=4pt,topsep=2pt]
  \item \emph{Independent convergence on the framework's structural choices.} The aggregate ratings on the four general-framing items, on the per-guideline ratings (twenty-six of thirty within plus or minus zero point four), and on the spectrum and gate mechanisms converge between the two groups, providing independent confirmation that the framework's structural choices are read consistently across two distinct engagement modes.
  \item \emph{Quantitative reinforcement of the soft-gates reformulation candidate.} The $\varepsilon_m$ calibration distribution shifts toward \emph{Too strict to be practically achievable} in the form-only group, adding independent quantitative weight to the soft-gates-to-hard-gates reformulation proposed in Interview 12.
  \item \emph{Quantitative reinforcement of the Organizational Alignment cross-axis reformulation.} The G7 Team-Architecture Balance ranking rises in the form-only group, with G9 and G12 falling, adding independent quantitative weight to the cross-axis reformulation of Organizational Alignment that surfaced in Interviews 9, 10, and 13.
  \item \emph{Quantitative reinforcement of the cost-conditioned G6 framing.} The G6 adoption likelihood and ranking drop in the form-only group, adding independent quantitative weight to the cost-conditioned and timing-based framings of G6 that surfaced in Interviews 12, 13, and 14.
  \item \emph{Independent endorsement of the deliberate-destination framing.} The form-only respondents endorse the modular monolith as a deliberate destination at substantively equivalent levels to the interview-and-questionnaire respondents, indicating the framing is not dependent on the conversational anchoring of an interview.
\end{enumerate}
